% NE 630 -- Lesson 23 handout
% Compile: latexmk -xelatex -interaction=nonstopmode -halt-on-error lesson_23_handout.tex
% This source also compiles with pdflatex when using the built-in fallback.
%
% Use the course's shared reading-guide style when available.
% The embedded fallback keeps the draft portable outside the repository.
% To force the fallback, replace the next line with \def\NEOriginalStyle{0}.
\IfFileExists{ne630_ch04_reading_guide_common.tex}{\def\NEOriginalStyle{1}}{\def\NEOriginalStyle{0}}
\ifnum\NEOriginalStyle=1\relax
  % NE 630 -- shared reading-guide definitions, Lessons 15--21
% Included by lesson_15_handout.tex through lesson_21_handout.tex.
% Layout/palette adapted from the user's ne630boardhandout class.
% Content: supplied chapter(1).tex and lecture_15.tex through lecture_19.tex.
% The supplied chapter numbers the four-factor discussion as Section 4.5.
% Drawings below adapt supplied Fig. 4.1, Fig. 4.5, Fig. 4.6, and the L19
% pin-cell drawing. They are qualitative schematics, not simulation results.
\documentclass[10pt,letterpaper]{article}
\usepackage{amsmath,mathtools}
\usepackage{libertinus,microtype}
\usepackage[left=0.48in,right=0.48in,top=0.48in,bottom=0.44in,
  headheight=20pt,headsep=5pt,footskip=17pt,includeheadfoot]{geometry}
\usepackage{graphicx,xcolor,tabularx,booktabs,array,ragged2e}
\usepackage{tikz}
\usetikzlibrary{arrows.meta,calc,positioning,patterns}
\usepackage[most]{tcolorbox}
\usepackage{fancyhdr,lastpage}
\usepackage[hidelinks]{hyperref}
\hypersetup{pdfauthor={NE 630},pdfsubject={The Power Reactor Core}}
\definecolor{KSUPurple}{HTML}{512888}
\definecolor{KSUDeepPurple}{HTML}{3B1B64}
\definecolor{KSUOrange}{HTML}{CA7C1B}
\definecolor{KSUGray}{HTML}{5D5D5D}
\setlength{\parindent}{0pt}
\setlength{\parskip}{3pt}
\setlength{\emergencystretch}{1.5em}
\setlength{\abovedisplayskip}{4pt}
\setlength{\belowdisplayskip}{4pt}
\setlength{\abovedisplayshortskip}{2pt}
\setlength{\belowdisplayshortskip}{3pt}
\setlength{\jot}{3pt}
\renewcommand{\arraystretch}{1.12}
\newcolumntype{Y}{>{\RaggedRight\arraybackslash}X}
\newcommand{\LessonNo}{15}
\pagestyle{fancy}
\fancyhf{}
\fancyhead[L]{\sffamily\bfseries\small NE 630}
\fancyhead[C]{\sffamily\small The Power Reactor Core}
\fancyhead[R]{\sffamily\bfseries\small Lesson \LessonNo}
\fancyfoot[L]{\sffamily\scriptsize  }
\fancyfoot[C]{\sffamily\scriptsize \thepage\ of \pageref*{LastPage}}
\fancyfoot[R]{\sffamily\scriptsize }
\renewcommand{\headrulewidth}{0.5pt}
\makeatletter
\renewcommand{\headrule}{\hbox to\headwidth{\color{KSUPurple}\leaders\hrule height .5pt\hfill}}
\makeatother
\newcommand{\Lesson}[4]{%
  \renewcommand{\LessonNo}{#1}%
  \noindent\begin{minipage}[c]{0.025\linewidth}\color{KSUPurple}\rule{3.5pt}{32pt}\end{minipage}%
  \begin{minipage}[c]{0.815\linewidth}
   {\sffamily\bfseries\fontsize{18}{20}\selectfont #2}\par
   {\sffamily\fontsize{9}{10.5}\selectfont #4}
  \end{minipage}\hfill
  \begin{minipage}[c]{0.14\linewidth}\raggedleft
    {\sffamily\color{KSUPurple}\bfseries\fontsize{25}{26}\selectfont #1}\par
    {\sffamily\scriptsize LESSON}
  \end{minipage}\par\vspace{3pt}
  \begin{tcolorbox}[enhanced,boxrule=.45pt,arc=1pt,colback=KSUPurple!4,
   colframe=KSUPurple!40,left=5pt,right=5pt,top=3pt,bottom=3pt,
   before skip=0pt,after skip=6pt]
   \sffamily\fontsize{9}{10.5}\selectfont\textbf{Read:} #3
  \end{tcolorbox}}
\newcommand{\Sec}[2]{\par\addvspace{5pt}\noindent
  \colorbox{KSUPurple}{\color{white}\sffamily\bfseries\fontsize{9}{10}\selectfont\strut #1}
  \hspace{3pt}{\sffamily\bfseries\color{KSUDeepPurple}\fontsize{10.2}{11.5}\selectfont #2}\par\vspace{1pt}}
\newcommand{\eqread}[1]{{\sffamily\fontsize{7.9}{9}\selectfont\color{KSUGray}#1}}
\newcommand{\figcap}[1]{\par\vspace{1pt}{\sffamily\fontsize{7.7}{9}\selectfont\color{KSUGray}#1}\par}
\newtcolorbox{KeyBox}[1]{enhanced,boxrule=.45pt,arc=1pt,
 colback=KSUPurple!3,colframe=KSUPurple!40,left=5pt,right=5pt,top=4pt,bottom=4pt,
 before skip=4pt,after skip=4pt,before upper=\RaggedRight,fonttitle=\sffamily\bfseries\small,
 coltitle=KSUDeepPurple,title={#1},colbacktitle=KSUPurple!7,
 titlerule=0pt,toptitle=2pt,bottomtitle=2pt}
\newtcolorbox{Explore}[1]{enhanced,boxrule=.5pt,arc=1pt,
 colback=KSUOrange!5,colframe=KSUOrange!65,left=6pt,right=6pt,top=4pt,bottom=4pt,
 before skip=7pt,after skip=3pt,before upper=\RaggedRight,
 title={OpenMC \enspace\textbar\enspace #1},fonttitle=\sffamily\bfseries\fontsize{9}{10}\selectfont,
 coltitle=KSUDeepPurple,colbacktitle=KSUOrange!10,titlerule=0pt,toptitle=2pt,bottomtitle=2pt}
\newcommand{\SourceLine}[1]{\par\vspace{2pt}{\sffamily\fontsize{7.2}{8.4}\selectfont\color{KSUGray}#1}\par}
\newcommand{\colstart}{\begin{minipage}[t]{.482\linewidth}\vspace{0pt}\RaggedRight}
\newcommand{\colnext}{\end{minipage}\hfill\begin{minipage}[t]{.482\linewidth}\vspace{0pt}\RaggedRight}
\newcommand{\colend}{\end{minipage}\par}
\newcommand{\uunit}[1]{\,\mathrm{#1}}
\newcommand{\metal}{\mathrm{metal}}
\newcommand{\water}{\mathrm{w}}
\newcommand{\fiiso}{\mathrm{fi}}
\newcommand{\feiso}{\mathrm{fe}}
\newcommand{\nucl}[2]{\ensuremath{{}^{#1}\mathrm{#2}}}



% Five reactor motifs, adapted from the supplied Fig. 4.1, with a compact TRISO zoom.
\newcommand{\FamilyFigure}{%
\begin{tikzpicture}[x=1cm,y=1cm,line width=.42pt,font=\sffamily\fontsize{7.5}{8.5}\selectfont,
 every node/.style={inner sep=1pt},>=Latex]
\begin{scope}[shift={(0,0)}]
 \draw[fill=black!3] (-.60,-.50) rectangle (.60,.50);
 \draw[fill=black!35] (0,0) circle (.34);
 \draw[fill=KSUPurple!27] (0,0) circle (.29);
 \node at (0,0) {Fuel}; \node at (0,.42) {H$_2$O};
 \node[align=center] at (0,-.84) {\textbf{LWR}\\square pin lattice};
\end{scope}
\begin{scope}[shift={(3.05,0)}]
 \draw[fill=black!3] (90:.67)--(30:.67)--(-30:.67)--(-90:.67)--(-150:.67)--(150:.67)--cycle;
 \draw[fill=black!35] (0,0) circle (.38);
 \draw[fill=KSUPurple!27] (0,0) circle (.32);
 \node at (0,0) {Fuel}; \node at (0,.50) {Na};
 \node[align=center] at (0,-.84) {\textbf{SFR}\\triangular pin lattice};
\end{scope}
\begin{scope}[shift={(6.10,0)}]
 \draw[fill=black!7] (-.78,-.50) rectangle (.78,.50);
 \draw[fill=white] (0,0) circle (.39);
 \draw (0,0) circle (.34);
 \foreach \ang in {0,60,...,300}{\draw[fill=KSUPurple!27] (\ang:.21) circle (.07);} 
 \draw[fill=KSUPurple!27] (0,0) circle (.07);
 \node[align=center] at (-.58,0) {D$_2$O};
 \node[align=center] at (0,-.84) {\textbf{PHWR / CANDU}\\bundle in a pressure tube};
\end{scope}
\begin{scope}[shift={(9.25,0)}]
 \draw[fill=black!12] (-.84,-.50) rectangle (.84,.50);
 \foreach \x/\y in {-.50/.25,0/.25,.50/.25,-.25/-.20,.25/-.20}{\draw[fill=white] (\x,\y) circle (.10);} 
 \foreach \x/\y in {-.56/-.18,-.19/.03,.19/.03,.56/-.18}{
   \draw[fill=KSUPurple!10,pattern=dots,pattern color=KSUPurple!70] (\x,\y) circle (.085);
   \draw (\x,\y) circle (.085);
 }
 \node[anchor=west] at (.86,.22) {He};
 \draw[->] (.82,.22)--(.59,.22);
 \node[anchor=west] at (.86,-.18) {TRISO};
 \draw[->] (.82,-.18)--(.62,-.18);
 \node[align=center] at (0,-.84) {\textbf{HTGR}\\TRISO in graphite};
\end{scope}
\begin{scope}[shift={(12.55,0)}]
 \draw[fill=KSUPurple!70] (0,0) circle (.10);  % kernel
 \draw[fill=KSUPurple!18] (0,0) circle (.18);  % buffer
 \draw[fill=gray!25] (0,0) circle (.24);       % inner PyC
 \draw[fill=green!12!black!8] (0,0) circle (.30); % SiC
 \draw[fill=gray!8] (0,0) circle (.36);        % outer PyC
 \foreach \r in {.10,.18,.24,.30,.36}{\draw (0,0) circle (\r);} 
 \node[font=\tiny] at (0,0) {K};
 \node[font=\tiny,anchor=west] at (.62,.25) {OPyC}; \draw[->] (.58,.24)--(.33,.24);
 \node[font=\tiny,anchor=west] at (.62,.06) {SiC};  \draw[->] (.58,.05)--(.28,.05);
 \node[font=\tiny,anchor=west] at (.62,-.12) {IPyC};\draw[->] (.58,-.11)--(.21,-.11);
 \node[font=\tiny,anchor=west] at (.62,-.30) {buffer};\draw[->] (.58,-.29)--(.14,-.20);
 \node[font=\tiny,anchor=east] at (-.62,.32) {kernel};\draw[->] (-.58,.30)--(-.08,.08);
 \node[align=center] at (0,-.84) {\textbf{TRISO particle}\\coated-fuel cross section};
\end{scope}
\end{tikzpicture}}

% Square and hexagonal pin cells, adapted from the L19 cell figure and Fig. 4.1.
\newcommand{\PinCellsFigure}{%
\begin{tikzpicture}[x=1cm,y=1cm,font=\sffamily\fontsize{8.5}{10}\selectfont,
 line width=.55pt,>=Latex,every node/.style={inner sep=1pt}]
\begin{scope}[shift={(0,0)}]
 \draw[fill=black!3] (-.88,-.88) rectangle (.88,.88);
 \draw[fill=KSUOrange!35] (0,0) circle (.58);
 \draw[fill=KSUPurple!25] (0,0) circle (.49);
 \draw[->] (0,0)--(45:.49) node[midway,above left=-1pt] {$r_f$};
 \draw[->] (0,0)--(-35:.58);\node[anchor=west] at (.57,-.40) {$r_c$};
 \draw[|<->|] (-.88,-1.10)--(.88,-1.10) node[midway,below=3pt] {$P$};
 \node at (0,1.10) {Square};
\end{scope}
\begin{scope}[shift={(3.80,0)}]
 \draw[fill=black!3] (90:1.016)--(30:1.016)--(-30:1.016)--(-90:1.016)--(-150:1.016)--(150:1.016)--cycle;
 \draw[fill=KSUOrange!35] (0,0) circle (.58);
 \draw[fill=KSUPurple!25] (0,0) circle (.49);
 \draw[densely dashed] (0,0)--(.88,0);
 \node[above] at (.69,.05) {$P/2$};
 \draw[|<->|] (-.88,-1.22)--(.88,-1.22) node[midway,below=3pt] {$P$};
 \node at (0,1.24) {Hexagonal};
 \node[anchor=west] at (.82,.77) {$a$};
\end{scope}
\end{tikzpicture}}

% Bookkeeping cycle adapted from supplied Fig. 4.5 (not a spatial plot).
\newcommand{\CycleFigure}{%
\begin{tikzpicture}[x=1cm,y=1cm,>=Latex,
 box/.style={draw=KSUPurple!45,fill=KSUPurple!3,rounded corners=1pt,
 minimum width=3.0cm,minimum height=1.05cm,align=center,inner sep=3pt},
 every node/.style={font=\sffamily\fontsize{8.2}{9.5}\selectfont}]
 \node[box] (n) at (0,0) {$n$\\thermal-fission births};
 \node[box] (en) at (3.85,0) {$\epsilon n$\\all fission births};
 \node[box] (epn) at (7.7,0) {$\epsilon p n$\\reach thermal energies};
 \node[box] (epfn) at (11.55,0) {$\epsilon p f n$\\absorbed in fuel};
 \node[box] (kn) at (15.4,0) {$\epsilon p f\eta_T n=k_\infty n$\\next thermal-fission births};
 \draw[->,line width=.6pt] (n)--(en) node[midway,above=2pt] {$\epsilon$};
 \draw[->,line width=.6pt] (en)--(epn) node[midway,above=2pt] {$p$};
 \draw[->,line width=.6pt] (epn)--(epfn) node[midway,above=2pt] {$f$};
 \draw[->,line width=.6pt] (epfn)--(kn) node[midway,above=2pt] {$\eta_T$};
\end{tikzpicture}}

% Qualitative Doppler/space-energy shielding sketch, adapted from Fig. 4.6.
\newcommand{\ResonanceFigure}{%
\begin{tikzpicture}[x=1cm,y=1cm,>=Latex,
 ax/.style={->,line width=.5pt,draw=black!75},
 cold/.style={line width=.85pt,draw=KSUPurple},
 hot/.style={line width=.85pt,draw=KSUOrange,dashed},
 every node/.style={font=\sffamily\fontsize{8.1}{9}\selectfont,inner sep=1pt}]
\begin{scope}
 \draw[ax] (0,0)--(3.5,0) node[right] {$E$};
 \draw[ax] (0,0)--(0,1.9) node[above] {$\sigma_a(E)$};
 \draw[cold] plot[smooth] coordinates {(0,.01)(.5,.04)(.9,.15)(1.2,.52)(1.4,1.2)(1.55,1.7)(1.7,1.2)(1.95,.43)(2.3,.10)(2.8,.03)(3.2,.01)};
 \draw[hot] plot[smooth] coordinates {(0,.02)(.5,.10)(.9,.28)(1.2,.72)(1.55,1.08)(1.9,.75)(2.3,.30)(2.8,.10)(3.2,.03)};
 \node[anchor=west,text=KSUPurple] at (.08,1.63) {$T_1$};
 \node[anchor=west,text=KSUOrange] at (2.26,1.09) {$T_2>T_1$};
\end{scope}
\begin{scope}[shift={(4.25,0)}]
 \draw[ax] (0,0)--(3.5,0) node[right] {$E$};
 \draw[ax] (0,0)--(0,1.9) node[above] {$\varphi_f/\varphi_m$};
 \draw[black!40,densely dotted] (0,1.5)--(3.2,1.5);
 \node[anchor=east] at (-.05,1.5) {$1$};
 \draw[cold] plot[smooth] coordinates {(0,1.5)(.6,1.49)(1,1.3)(1.25,.65)(1.5,.3)(1.75,.59)(2,1.26)(2.5,1.49)(3.2,1.5)};
 \draw[hot] plot[smooth] coordinates {(0,1.5)(.4,1.43)(.8,1.22)(1.2,.88)(1.55,.65)(1.9,.83)(2.3,1.18)(2.8,1.43)(3.2,1.49)};
\end{scope}
\end{tikzpicture}}

\else
  \documentclass[10pt,letterpaper]{article}
  \usepackage[margin=0.50in,top=0.47in,bottom=0.48in,footskip=0.23in]{geometry}
  \usepackage[T1]{fontenc}
  \usepackage{newtxtext,newtxmath}
  \usepackage{microtype}
  \usepackage{amsmath,booktabs,array,graphicx,xcolor}
  \usepackage{ragged2e}
  \usepackage{tikz}
  \usepackage[most]{tcolorbox}
  \usepackage{fancyhdr}
  \usepackage[hidelinks]{hyperref}
  \definecolor{ksupurple}{HTML}{512888}
  \definecolor{orange}{HTML}{CA7C1B}
  \setlength{\parindent}{0pt}
  \setlength{\parskip}{2pt}
  \setlength{\emergencystretch}{1em}
  \pagestyle{fancy}
  \fancyhf{}
  \renewcommand{\headrulewidth}{0pt}
  \renewcommand{\footrulewidth}{0.3pt}
  \fancyfoot[L]{\sffamily\fontsize{8}{9}\selectfont NE 630\enspace /\enspace Nuclear Reactor Theory}
  \fancyfoot[R]{\sffamily\fontsize{8}{9}\selectfont Lesson 23\enspace /\enspace \thepage}
  \newcommand{\Lesson}[4]{%
    {\sffamily\bfseries\fontsize{9}{11}\selectfont\color{ksupurple}NE 630\enspace |\enspace LESSON #1}\par
    \vspace{1pt}
    {\sffamily\bfseries\fontsize{19}{22}\selectfont #2\par}
    \vspace{2pt}
    {\fontsize{8.5}{10}\selectfont\textbf{Reading:} #3\par}
    {\fontsize{10}{12}\selectfont\color{ksupurple}\textit{#4}\par}
    \vspace{2pt}{\color{ksupurple}\hrule height 0.7pt}\vspace{5pt}}
  \newcommand{\Sec}[2]{\par\addvspace{6pt}%
    {\sffamily\bfseries\color{ksupurple}#1\enspace #2\par}\vspace{2pt}}
  \newcommand{\colstart}{\par\noindent\begin{minipage}[t]{0.483\textwidth}\vspace{0pt}\RaggedRight}
  \newcommand{\colnext}{\end{minipage}\hfill\begin{minipage}[t]{0.483\textwidth}\vspace{0pt}\RaggedRight}
  \newcommand{\colend}{\end{minipage}\par}
  \newtcolorbox{KeyBox}[1]{enhanced,boxrule=0.45pt,arc=1mm,
    colback=ksupurple!2,colframe=ksupurple!65,
    colbacktitle=ksupurple!9,coltitle=ksupurple,
    fonttitle=\sffamily\bfseries\fontsize{9}{10}\selectfont,title={#1},
    left=5pt,right=5pt,top=3pt,bottom=3pt,
    toptitle=3pt,bottomtitle=3pt,before skip=5pt,after skip=5pt}
  \newtcolorbox{Explore}[1]{enhanced,boxrule=0.45pt,arc=1mm,
    colback=orange!4,colframe=orange!75,
    colbacktitle=orange!12,coltitle=black,
    fonttitle=\sffamily\bfseries\fontsize{9}{10}\selectfont,title={#1},
    left=6pt,right=6pt,top=3pt,bottom=3pt,
    toptitle=3pt,bottomtitle=3pt,before skip=7pt,after skip=0pt}
  \newcommand{\figcap}[1]{{\fontsize{8}{9.5}\selectfont #1\par}}
\fi

% Local layout helpers; shared definitions are unchanged for other lessons.
\usepackage{amsmath,booktabs,array,graphicx,xcolor,tikz,ragged2e}
\usetikzlibrary{arrows.meta,calc,decorations.pathreplacing,positioning}
\providecolor{ksupurple}{HTML}{512888}
% Match the Wildcat theme's wcalerted, overriding xcolor's built-in orange.
\definecolor{orange}{HTML}{CA7C1B}
% These synthesis prompts are not OpenMC-specific.
\renewtcolorbox{Explore}[1]{enhanced,boxrule=.5pt,arc=1pt,
 colback=orange!5,colframe=orange!65,left=6pt,right=6pt,top=4pt,bottom=4pt,
 before skip=6pt,after skip=3pt,before upper=\RaggedRight,
 title={#1},fonttitle=\sffamily\bfseries\fontsize{9}{10}\selectfont,
 coltitle=ksupurple,colbacktitle=orange!10,titlerule=0pt,
 toptitle=2pt,bottomtitle=2pt}
\newcommand{\NEContinue}[2]{%
  {\sffamily\bfseries\fontsize{12}{14}\selectfont\color{ksupurple}Lesson #1\enspace |\enspace #2\par}
  \vspace{4pt}{\color{ksupurple}\hrule height 0.6pt}\vspace{5pt}}
\newcommand{\NEEx}[1]{\par\addvspace{4pt}\textbf{#1}\par}
\newcommand{\NECheck}{\textbf{Definition check.}\ }
\newcommand{\NEsmall}[1]{{\fontsize{8.3}{10}\selectfont\RaggedRight #1\par}}
\newcommand{\NEbody}{%
  \fontsize{9}{10.7}\selectfont\RaggedRight
  \setlength{\parskip}{2pt}
  \setlength{\topsep}{2pt}\setlength{\partopsep}{0pt}
  \setlength{\abovedisplayskip}{4pt}\setlength{\belowdisplayskip}{4pt}
  \setlength{\abovedisplayshortskip}{3pt}\setlength{\belowdisplayshortskip}{3pt}}


% FNRP Chapter 10; both numerical examples retain their supplied data.
% The editable chain is included directly from figures/depletion_chain_tikz.tex.
% No fuel density or enrichment convention is silently supplied.
\hypersetup{pdftitle={NE 630 -- Lesson 23: Reactivity impacts of fuel depletion},pdfauthor={Jeremy Roberts},pdfsubject={Fuel depletion and long-term core reactivity}}
\begin{document}
\NEbody
\Lesson{23}{Reactivity impacts of fuel depletion}
 {FNRP \S\,10.4--10.5; Eqs.\ (10.27)--(10.39).}
 {}
\textbf{Objective:} Quantify the evolution of nuclides and core reactivity with burnup.

\colstart
\Sec{1}{Review: reactivity in time}
As a reactor operates (at power), its excess reactivity typically decreases
because \textcolor{ksupurple}{${}^{235}$U is consumed} and
\textcolor{ksupurple}{parasitic fission products (FPs) are born}. However,
that excess reactivity must \textcolor{orange}{be balanced by control}.

\vspace{5pt}

To represent these processes, define a time-varying multiplication factor as
\begin{equation*}
 k(t)=\frac{\textcolor{ksupurple}{\nu\bar\Sigma_{fT}^{F}(t)}
                 \,p\,\varepsilon P_{\mathrm{NL}}}
 {\bar\Sigma_{aT}^{F}(t)
  +\textcolor{ksupurple}{\bar\Sigma_{aT}^{\mathrm{FP}}(t)}
  +\varsigma\frac{V_M}{V_F}
   \left(\bar\Sigma_{aT}^{M}
   +\textcolor{orange}{\bar\Sigma_{aT}^{\mathrm{control}}(t)}\right)},
 \tag{FNRP 10.4*}
\end{equation*}
where $P_{\mathrm{NL}}$ is the probability that neutrons do not leak from the system.
The asterisk marks the explicit control term.



\vspace{5pt}

During operation, $\textcolor{orange}{\bar\Sigma_{aT}^{\mathrm{control}}(t)}$
must be adjusted to maintain $k(t)=1$. Here,
\textcolor{orange}{adjustable control} is introduced within the moderator
volume, consistent with typical LWRs, whether by changing the
\textcolor{orange}{rod position} or \textcolor{orange}{soluble poison concentration}.

\vspace{5pt}

\Sec{2}{A representative depletion chain}

In theory, all nuclides connected through capture or decay should be
simultaneously modeled as functions of time. A practical depletion chain for
analysis of LWRs is shown below: each arrow \emph{into} a nuclide represents
production, while an \emph{outward} arrow represents destruction.

\begin{center}
\begingroup
\def\NEDepletionChainEmbedded{}
% Editable, column-width adaptation of the supplied depletion_chain_tikz.tex.
% Standalone: xelatex depletion_chain_tikz.tex
% Embedded: define \NEDepletionChainEmbedded, then \input this file in a group.
% Original Downloads file is retained unchanged. The source palette is retained.
% Routine route numbers are omitted; the handout caption explains the key.
% Reference: https://www-nds.iaea.org/wimsd/achain.htm
% Horizontal coordinates are shared mass-number columns (x = A - 231).
% Restore the source's cross-row alignments; never center rows independently.
% Scale geometry only: fixed-size text is not transformed or resized.
\ifdefined\NEDepletionChainEmbedded
  \let\NEChainStart\relax
  \let\NEChainFinish\relax
\else
  \documentclass[tikz,border=3pt]{standalone}
  \usepackage{libertinus}
  \usepackage{amsmath}
  \newcommand{\NEChainStart}{\begin{document}}
  \newcommand{\NEChainFinish}{\end{document}}
\fi
\NEChainStart
\begingroup
\usetikzlibrary{arrows.meta,calc,backgrounds}
\definecolor{NEChainPurple}{HTML}{512888}
\definecolor{NEChainOrange}{HTML}{F58220}
\definecolor{NEChainLabel}{HTML}{202024}
% Compact chart labels keep the isotope centers on their physical columns.
% Mass/state above element avoids squeezing adjacent labels horizontally.
\newcommand{\NEChainNuclide}[2]{\shortstack[c]{%
  {\fontsize{5.8}{6.2}\selectfont\ensuremath{#1}}\\[-1.5pt]
  {\fontsize{7.4}{8.5}\selectfont\ensuremath{\mathrm{#2}}}}}
\begin{tikzpicture}[
  x=.655cm,y=.64cm,
  font=\fontsize{7.4}{8.5}\selectfont,
  nuclide/.style={text=NEChainLabel,fill=white,
    inner xsep=.35pt,inner ysep=1pt,outer sep=.2pt},
  pseudo/.style={text=NEChainOrange,font=\fontsize{6.8}{8}\selectfont\itshape,
    inner sep=.2pt,anchor=north},
  orange line/.style={draw=NEChainOrange,line width=.45pt},
  purple line/.style={draw=NEChainPurple,line width=.45pt},
  orange arrow/.style={orange line,-{Stealth[length=1mm,width=.7mm]}},
  purple arrow/.style={purple line,-{Stealth[length=1mm,width=.7mm]}}
]
% ------------------------------------------------------------------
% Nuclides. Equal A has equal x across every element row.
% Element rows retain the supplied vertical coordinates.
% ------------------------------------------------------------------
\node[nuclide] (Th232) at (1.0000,0.0000) {\NEChainNuclide{232}{Th}};

\node[nuclide] (Pa231) at (0.0000,1.5000) {\NEChainNuclide{231}{Pa}};
\node[nuclide] (Pa233) at (2.0000,1.5000) {\NEChainNuclide{233}{Pa}};

\node[nuclide] (U232)  at (1.0000,3.0000) {\NEChainNuclide{232}{U}};
\node[nuclide] (U233)  at (2.0000,3.0000) {\NEChainNuclide{233}{U}};
\node[nuclide] (U234)  at (3.0000,3.0000) {\NEChainNuclide{234}{U}};
\node[nuclide] (U235)  at (4.0000,3.0000) {\NEChainNuclide{235}{U}};
\node[nuclide] (U236)  at (5.0000,3.0000) {\NEChainNuclide{236}{U}};
\node[nuclide] (U237)  at (6.0000,3.0000) {\NEChainNuclide{237}{U}};
\node[nuclide] (U238)  at (7.0000,3.0000) {\NEChainNuclide{238}{U}};

\node[nuclide] (Np237) at (6.0000,4.6500) {\NEChainNuclide{237}{Np}};
\node[nuclide] (Np239) at (8.0000,4.6500) {\NEChainNuclide{239}{Np}};

\node[nuclide] (Pu238) at (7.0000,6.3500) {\NEChainNuclide{238}{Pu}};
\node[nuclide] (Pu239) at (8.0000,6.3500) {\NEChainNuclide{239}{Pu}};
\node[nuclide] (Pu240) at (9.0000,6.3500) {\NEChainNuclide{240}{Pu}};
\node[nuclide] (Pu241) at (10.0000,6.3500) {\NEChainNuclide{241}{Pu}};
\node[nuclide] (Pu242) at (11.0000,6.3500) {\NEChainNuclide{242}{Pu}};

\node[nuclide] (Am241)  at (10.0000,7.9000) {\NEChainNuclide{241}{Am}};
\node[nuclide] (Am242m) at (11.0000,7.9000) {\NEChainNuclide{242\mathrm{m}}{Am}};
\node[nuclide] (Am243)  at (12.0000,7.9000) {\NEChainNuclide{243}{Am}};
\node[nuclide,text=NEChainOrange] (Am242g) at (11.0000,9.6500)
  {\NEChainNuclide{242\mathrm{g}}{Am}};

\node[nuclide] (Cm242) at (11.0000,11.2500) {\NEChainNuclide{242}{Cm}};
\node[nuclide] (Cm243) at (12.0000,11.2500) {\NEChainNuclide{243}{Cm}};
\node[nuclide] (Cm244) at (13.0000,11.2500) {\NEChainNuclide{244}{Cm}};

% The pseudo states remain unboxed labels, as in the source.
\coordinate (PseudoLeft)  at (1.5000,2.1500);
\coordinate (PseudoRight) at (6.5000,2.1500);
\node[pseudo] at (.3000,2.3000) {Pseudo};
\node[pseudo] at (6.5000,2.0500) {Pseudo};

% Routing points for two-color paths. The color changes occur at
% these bends; only the final segment carries an arrowhead.
\coordinate (Pa231Bend) at (0.0000,0.0000);
\coordinate (Pa233Bend) at (2.0000,0.0000);
\coordinate (U232Bend) at (1.0000,1.5000);
\coordinate (U234Bend) at (3.0000,1.5000);
\coordinate (Np239Bend) at (8.0000,3.0000);
\coordinate (Pu238Bend) at (7.0000,4.6500);
\coordinate (Pu240Bend) at (9.0000,4.9700);
\coordinate (Am243Bend) at (12.0000,6.3500);
\coordinate (Cm244Bend) at (13.0000,7.9000);

% Draw paths behind the labels. Crossings without a dot are retained
% as ordinary crossings, not extra graph nodes.
\begin{pgfonlayer}{background}

% ------------------------------------------------------------------
% Branch notes below refer to the IAEA web page, not visible figure labels.
% ------------------------------------------------------------------
% (1), (2): horizontal curium paths.
\draw[orange arrow] (Cm243.east) -- (Cm244.west);
\draw[orange arrow] (Cm242.east) -- (Cm243.west);

% (3), (4): vertically stacked americium states and curium.
\draw[purple arrow] (Am242g.north) -- (Cm242.south);
\draw[purple arrow] (Am242m.north) -- (Am242g.south);

% (5), (6): the two branches out of 241Am.
\draw[orange arrow] (Am241.east) -- (Am242m.west);
\draw[orange arrow] (Am241.north east) -- (Am242g.south west);

% (7): leftward return from 241Am to 237Np.
\draw[purple arrow] (Am241.west) -- (6.0000,7.9000) -- (Np237.north);

% (8)--(10): three nested returns from Cm to Pu.
\draw[purple arrow] (Cm244.north) -- (13.0000,13.0500) --
  (9.0000,13.0500) -- (Pu240.north);
\draw[purple arrow] (Cm243.north) -- (12.0000,12.5000) --
  (8.0000,12.5000) -- (Pu239.north);
\draw[purple arrow] (Cm242.north) -- (11.0000,11.9500) --
  (7.0000,11.9500) -- (Pu238.north);

% (11): 241Pu to 241Am.
\draw[purple arrow] (Pu241.north) -- (Am241.south);

% (12): rightward orange segment, then upward purple segment.
\draw[orange line] (Np239.north east) -- (Np239.north east |- Pu240Bend)
  -- (Pu240Bend);
\draw[purple arrow] (Pu240Bend) -- (Pu240.south);

% (13): direct vertical path from 239Np to 239Pu.
\draw[purple arrow] (Np239.north) -- (Pu239.south);

% (14): 237Np to 238Pu, with the original color change at the bend.
\draw[orange line] (Np237.east) -- (Pu238Bend);
\draw[purple arrow] (Pu238Bend) -- (Pu238.south);

% (15): leftward/downward return from 238Pu to 234U.
\draw[purple arrow] (Pu238.west) -- (3.0000,6.3500) -- (U234.north);

% (16): 238U to 239Np, with a color change at the bend.
\draw[orange line] (U238.east) -- (Np239Bend);
\draw[purple arrow] (Np239Bend) -- (Np239.south);

% IAEA (19): 237U to 237Np (the downloaded drawing swapped 17/19).
\draw[purple arrow] (U237.north) -- (Np237.south);

% (18): 233Pa to 234U, with a color change at the bend.
\draw[orange line] (Pa233.east) -- (U234Bend);
\draw[purple arrow] (U234Bend) -- (U234.south);

% IAEA (17): right-hand pseudo route, 238U -> Pseudo -> 237U.
\draw[orange arrow] (U238.south west) -- (PseudoRight);
\draw[purple arrow] (PseudoRight) -- (U237.south east);

% (20): the final horizontal U-chain step.
\draw[orange arrow] (U237.east) -- (U238.west);

% (21): the left-hand pseudo route, 233U -> Pseudo -> 232U.
\draw[orange arrow] (U233.south west) -- (PseudoLeft);
\draw[purple arrow] (PseudoLeft) -- (U232.south east);

% (22), (23): the two paths out of 232Th.
\draw[orange line] (Th232.west) -- (Pa231Bend);
\draw[purple arrow] (Pa231Bend) -- (Pa231.south);
\draw[orange line] (Th232.east) -- (Pa233Bend);
\draw[purple arrow] (Pa233Bend) -- (Pa233.south);

% ------------------------------------------------------------------
% Paths that are unnumbered in the supplied figure.
% ------------------------------------------------------------------
% Pa -> U paths.
\draw[orange line] (Pa231.east) -- (U232Bend);
\draw[purple arrow] (U232Bend) -- (U232.south);
\draw[purple arrow] (Pa233.north) -- (U233.south);

% Main uranium chain (the last step, (20), is drawn above).
\foreach \from/\to in {U233/U234,U234/U235,U235/U236,U236/U237}
  \draw[orange arrow] (\from.east) -- (\to.west);

% Main plutonium chain.
\foreach \from/\to in {Pu238/Pu239,Pu239/Pu240,Pu240/Pu241,Pu241/Pu242}
  \draw[orange arrow] (\from.east) -- (\to.west);

% Both americium states connect to 243Am.
\draw[orange arrow] (Am242m.east) -- (Am243.west);
\draw[orange arrow] (Am242g.south east) -- (Am243.north west);

% Capture in 242Pu, followed by decay of the omitted intermediate, feeds 243Am.
\draw[orange line] (Pu242.east) -- (Am243Bend);
\draw[purple arrow] (Am243Bend) -- (Am243.south);

% Right-hand path from 243Am to 244Cm.
\draw[orange line] (Am243.east) -- (Cm244Bend);
\draw[purple arrow] (Cm244Bend) -- (Cm244.south);

% Outer return from 243Am to 239Np, below the Pu row.
\draw[purple arrow] (Am243.south east) -- (12.2045,7.5500) --
  (12.2045,4.6500) -- (Np239.east);

\end{pgfonlayer}
\end{tikzpicture}
\endgroup
\NEChainFinish

\endgroup
\end{center}
\figcap{Adapted from the \href{https://www-nds.iaea.org/wimsd/achain.htm}{IAEA WIMS-D chain}.
\textcolor[HTML]{F58220}{Orange}: neutron-induced production;
\textcolor[HTML]{512888}{purple}: decay.
Mixed (orange+purple) paths collapse immediately decaying intermediates. The superscripts
$\mathrm{m}$ and $\mathrm{g}$ refer to metastable and ground states; orange ${}^{242\mathrm{g}}$Am
uses reduced-concentration bookkeeping for its branch to ${}^{242}$Cm.
\textit{Pseudo} denotes a fictitious fission product that is introduced to account for non-negligible contributions from $(n,2n)$ interactions.  See the IAEA WIMS-D documentation for details. }

\colnext

\Sec{3}{Further simplifications}

Because each fissile nuclide has unique fission-product distributions,
tracking the entire isotopic inventory remains quite complex!
A much more drastic simplification is presented in FNRP based on these assumptions:

\vspace{5pt}

\textbf{1.} Only ${}^{235}$U, ${}^{238}$U, and ${}^{239}$Pu are important.
\vspace{2pt}


\textbf{2.} Ignore their decay.
\vspace{2pt}

\textbf{3.} ${}^{239}$Pu is produced \emph{directly} and \emph{immediately}
from ${}^{238}$U$(n,\gamma)$ reactions.
\vspace{2pt}

\textbf{4.} The amount of ${}^{238}$U is so large that absorption losses
are negligible, and the amount of ${}^{238}$U remains constant.

\vspace{5pt}

The resulting model can then be expressed by the following three boxed equations:

\begin{equation*}
\tag{10.27 $\rightarrow$ 10.31}
\begin{split}
\frac{dN^{\text{U}235}}{dt} &=  -N^{\text{U}235}(t) \bar{\sigma}^{\text{U}235}_a \phi \\
 \longrightarrow \, & \boxed{N^{\text{U}235}(t) = N^{\text{U}235}(0) e^{-\bar{\sigma}^{\text{U}235}_a \phi t}}
\end{split}
\end{equation*}

\begin{equation*}
\tag{10.27 $\rightarrow$  10.34}
\begin{split}
\frac{dN^{\text{U}238}}{dt} &= -N^{\text{U}238}(t) \bar{\sigma}^{\text{U}238}_a \phi \approx 0 \\
 \longrightarrow \, & \boxed{N^{\text{U}238}(t) = N^{\text{U}238}(0)}
\end{split}
\end{equation*}

\begin{equation*}
\tag{10.30 $\rightarrow$  10.37}
\begin{split}
\frac{dN^{\text{Pu}239}}{dt} &=
 N^{\text{U}238}(t) \bar{\sigma}^{\text{U}238}_{\gamma} \phi
 - N^{\text{Pu}239}(t) \bar{\sigma}^{\text{Pu}239}_a \phi \\
\longrightarrow \, & \boxed{N^{\text{Pu}239}(t)
  = \frac{\bar{\sigma}^{\text{U}238}_{\gamma}}{\bar{\sigma}^{\text{Pu}239}_a}
    N^{\text{U}238}(0) \left [1 - e^{-\bar{\sigma}^{\text{Pu}239}_a \phi t} \right ]}
\end{split}
\end{equation*}


% For constant total flux $\phi$ and constant, spectrum-averaged cross sections,
% \begin{align*}
%  \frac{dN^{\mathrm{U}235}}{dt}
%  &=-N^{\mathrm{U}235}(t)\bar\sigma_a^{\mathrm{U}235}\phi,
%  \tag{FNRP 10.27}\\
%  N^{\mathrm{U}235}(t)
%  &=N^{\mathrm{U}235}(0)e^{-\bar\sigma_a^{\mathrm{U}235}\phi t}.
%  \tag{FNRP 10.31}
% \end{align*}
% Neglecting absorption losses in ${}^{238}$U,
% \begin{align*}
%  \frac{dN^{\mathrm{U}238}}{dt}
%  &=-N^{\mathrm{U}238}(t)\bar\sigma_a^{\mathrm{U}238}\phi\approx0,
%  \tag{FNRP 10.28}\\
%  N^{\mathrm{U}238}(t)&\simeq N^{\mathrm{U}238}(0).
%  \tag{FNRP 10.34}
% \end{align*}
% For ${}^{239}$Pu,
% \begin{equation*}
%  \frac{dN^{\mathrm{Pu}239}}{dt}
%  =N^{\mathrm{U}238}(0)\bar\sigma_\gamma^{\mathrm{U}238}\phi
%   -N^{\mathrm{Pu}239}(t)\bar\sigma_a^{\mathrm{Pu}239}\phi.
%  \tag{FNRP 10.35}
% \end{equation*}
% With $N^{\mathrm{Pu}239}(0)=0$, this gives
% \begin{equation*}
%  N^{\mathrm{Pu}239}(t)
%  =\frac{\bar\sigma_\gamma^{\mathrm{U}238}}{\bar\sigma_a^{\mathrm{Pu}239}}
%  N^{\mathrm{U}238}(0)\left[1-e^{-\bar\sigma_a^{\mathrm{Pu}239}\phi t}\right].
%  \tag{FNRP 10.37}
% \end{equation*}

\Sec{4}{Time, fluence, and burnup}
Time is a familiar quantity, but alternatives can be useful when quantifying
depletion effects. For example, the \textbf{fluence} is the time-integrated flux,
\begin{equation*}
 \Phi(t)=\int_0^t\phi(t')\,dt',
 \tag{FNRP 10.32}
\end{equation*}
and has units of $\mathrm{cm}^{-2}$. For constant flux, $\Phi(t)=\phi t$.

\vspace{5pt}

\textbf{Burnup} is the energy produced per (initial) unit mass of actinides
and is listed in a variety of units that include
\begin{itemize}\setlength{\itemsep}{1pt}\setlength{\parsep}{0pt}
 \item MWd/kg (most common),
 \item GWd/MTU (equivalent to MWd/kg on the same mass basis; MTU is
 ``metric ton (of) uranium''),
 \item J/kg, and
 \item \% ${}^{235}$U (not burnup, \emph{per se}, but what we use at our TRIGA).
\end{itemize}
\colend

\newpage

\colstart

\Sec{5}{Computing burnup with the simplified model}
The energy produced per unit volume in one second is the product of the
fission-rate density $R_f=\Sigma_f\phi$ [fissions/(cm$^3$s)] and the energy
released per fission, $\kappa\approx3.2\times10^{-11}$ J/fission
(assuming the usual 200 MeV/fission). For constant flux $\phi$, the
accumulated energy per unit volume over $t'\in[0,t]$ is
\begin{align*}
 Q'''(t) &= \int\int_0^t q'''(t')\,dt' \\
         &=\int_0^t\kappa\phi\,
 \overbrace{\Big(\bar\sigma_f^{\mathrm{U}235}N^{\mathrm{U}235}(t')
 +\bar\sigma_f^{\mathrm{U}238}N^{\mathrm{U}238}(t')
 +\bar\sigma_f^{\mathrm{Pu}239}N^{\mathrm{Pu}239}(t')\Big)}^{\Sigma_f(t')}\,dt' \, ,
\end{align*}
where $q'''(t)$ is the instantaneous power density.

\vspace{5pt}

The initial actinide mass per unit volume is
$m_0 \approx\rho_{\mathrm{UO_2}}(238/270)$, though enrichment should generally
be considered, and the corresponding burnup is
\[
 B(t)=\frac{Q'''(t)}{m_0}.
\]

\Sec{6}{Reactivity and depletion}

Ignoring control and fission products, the multiplication factor can now be written as
\begin{equation*}
 k(t)=\frac{\big(\nu\bar\sigma_{fT}^{\mathrm{U}235}N^{\mathrm{U}235}(t)
       +\nu\bar\sigma_{fT}^{\mathrm{Pu}239}N^{\mathrm{Pu}239}(t)\big)
       \varepsilon pP_{\mathrm{NL}}}
 {\bar\sigma_{aT}^{\mathrm{U}235}N^{\mathrm{U}235}(t)
  +\bar\sigma_{aT}^{\mathrm{Pu}239}N^{\mathrm{Pu}239}(t)
  +\Sigma_{aT}^{\mathrm{U}238}+\varsigma\frac{V_m}{V_f}\Sigma_{aT}^{m}}.
 \tag{FNRP 10.39*}
\end{equation*}

The $T$ subscripts make the thermal-group cross sections explicit.
Make sure to use the appropriate $\nu$ for each fissile nuclide.

\vspace{5pt}

\textbf{Important:} When doing depletion, we use the total flux over the
entire energy range and must use the appropriate \emph{1-group} cross sections.
When using the four-factor formula, we need the appropriate
\emph{thermal-group} cross sections!

\colnext
\Sec{7}{Example -- Part 1}
Consider 4\% enriched UO$_2$ fuel for which the relevant, spectrum-averaged
cross sections are

\vspace{5pt}

\begin{center}
\renewcommand{\arraystretch}{1.12}
\begin{tabular}{@{}lrr@{}}
\toprule
Nuclide & $\bar\sigma_a$ [b] & $\bar\sigma_f$ [b]\\
\midrule
${}^{235}$U  & $57.20$  & $46.89$\\
${}^{238}$U  & $1.02$   & $0.11$\\
${}^{239}$Pu & $195.10$ & $124.41$\\
\bottomrule
\end{tabular}
\end{center}


\vspace{5pt}

Assume the total flux is $\phi=2\times10^{14}$ cm$^{-2}$s$^{-1}$. Then
\begin{enumerate}\setlength{\itemsep}{2pt}\setlength{\parsep}{0pt}
 \item determine the number densities of ${}^{235}$U and ${}^{239}$Pu
 after each of 5 years of operation, and
 \item determine the fraction of all power generated due to ${}^{239}$Pu
 after 1 year of operation.
\end{enumerate}
\par\vspace{8pt}

\Sec{8}{Example -- Part 2}
Consider the depletion we did above for Example -- Part 1. The appropriate
\emph{thermal} cross sections for the fuel nuclides are

\vspace{5pt}

\begin{center}
\renewcommand{\arraystretch}{1.12}
\begin{tabular}{@{}lrrr@{}}
\toprule
Nuclide & $\bar\sigma_{aT}$ [b] & $\bar\sigma_{fT}$ [b] & $\bar\nu_T$\\
\midrule
${}^{235}$U  & $297.40$  & $253.09$ & $2.43$\\
${}^{238}$U  & $1.32$    & $0.00$   & $0.00$\\
${}^{239}$Pu & $1146.51$ & $732.22$ & $2.90$\\
\bottomrule
\end{tabular}
\end{center}

\vspace{5pt}

Assuming $\varsigma=1.1$, $V_m/V_f=1$, $p=0.6$, $\varepsilon=1.3$,
$P_{\mathrm{NL}}=0.95$, and $\Sigma_{aT}^{m}=0.007$ cm$^{-1}$, compute
\begin{enumerate}\setlength{\itemsep}{2pt}\setlength{\parsep}{0pt}
 \item the initial reactivity, and
 \item the reactivity, fluence, and burnup after each of 5 years of operation.
\end{enumerate}

\vspace{5pt}

\textbf{Ponderable:} How does the \emph{fission rate} change in time?

\colend
\end{document}
